% ===============================

% SEZIONE 1: Potenze e proprietà

% ===============================

\begin{question}[ID=1.1, subtitle={Semplificazione di espressioni con potenze}, tags={potenze, semplificazione}, difficulty=facile]
Semplificare le seguenti espressioni usando le proprietà delle potenze e scrivere il risultato sotto forma di potenza con esponente intero, quando possibile:
\begin{multicols}{2}
\begin{enumerate}[leftmargin=3.2cm]
    \item $2^{5}\cdot 2^{-3} : 2^{2}$
    \item $(3^{2}\cdot 3^{-1})^2$
    \item $\Big(\dfrac{2}{3}\Big)^3 : \Big(\dfrac{2}{3}\Big)^{-1} \cdot\Big(\dfrac{3}{2}\Big)^{-2} \cdot \Big(\dfrac{2}{3}\Big)^{-4}$
    \item $(5^{-1})^{-3}\cdot 5^{2}$
    \item $(2^{-3}\cdot 2^{5})^2 : 2^{2}$
    \item $(4^{3})^{-1}\cdot 4^{4}$
\end{enumerate}
\end{multicols}
\begin{solution}
\begin{enumerate}
    \item $2^0=1$
    \item $3^2=9$
    \item $\left(\tfrac{2}{3}\right)^2=\tfrac{4}{9}$
    \item $5^5=3125$
    \item $2^2=4$
    \item $4^1=4$
\end{enumerate}
\end{solution}
\end{question}

%----------------------------------------------------%

\begin{question}[ID=1.2, subtitle={Potenze con esponenti frazionari}, tags={potenze, esponenti frazionari}, difficulty=media]
Calcolare i seguenti valori e semplificare il più possibile:
\begin{multicols}{2}
\begin{enumerate}[leftmargin=3.2cm]
    \item $16^{3/4}$
    \item $27^{-2/3}$
    \item $\sqrt[3]{8^5}$
    \item Se $a>0$, esprimere $(a^{1/2}\cdot a^{-3/4})^{-2}$ come potenza di $a$ con esponente razionale.
    \item $81^{3/4}$
    \item Se $b>0$, esprimere $(b^{-1/2}\cdot b^{5/6})^3$ come potenza di $b$.
\end{enumerate}
\end{multicols}
\begin{solution}
\begin{enumerate}
    \item $8$
    \item $\tfrac{1}{9}$
    \item $8^{5/3}=32$
    \item $a^{1/2}$
    \item $27$
    \item $b$
\end{enumerate}
\end{solution}
\end{question}

%----------------------------------------------------%

\begin{question}[ID=1.3, subtitle={Equazioni risolte con le proprietà delle potenze}, tags={potenze, equazioni}, difficulty=media]
Risolvi le seguenti equazioni:
\begin{multicols}{2}
\begin{enumerate}[leftmargin=3.2cm]
    \item $4^{x}=64$
    \item $9\cdot 3^{2x-1}=81$
    \item $2^{2x+1}= \dfrac{1}{8}$
    \item $(x^{3})^{2}= 32$
    \item $5^{x+1}=125$
    \item $27^{x-1}=9$
\end{enumerate}
\end{multicols}
\begin{solution}
\begin{enumerate}
    \item $x=3$
    \item $x=\tfrac{3}{2}$
    \item $x=-2$
    \item $x=\pm2^{5/6}$
    \item $x=2$
    \item $x=\tfrac{5}{3}$
\end{enumerate}
\end{solution}
\end{question}

%----------------------------------------------------%

\begin{question}[ID=1.4, subtitle={Potenze di 10 e notazione scientifica}, tags={potenze, notazione scientifica}, difficulty=facile]
Potenze di 10 e notazione scientifica:
\begin{enumerate}[leftmargin=3.2cm]
    \item Riscrivi i seguenti numeri in notazione scientifica: $0{,}00052,\;1240000$.
    \item Calcola e semplifica: $(3{,}2\times 10^{4})\cdot (5\times 10^{-2})$.
    \item Semplifica: $\dfrac{10^{7}\cdot 10^{-3}}{10^{2}}$.
    \item Un valore in fisica è $6{,}02\times 10^{23}$. Scrivilo come potenza di 10 moltiplicata per un numero compreso tra 1 e 10 e poi calcola $(6{,}02\times 10^{23})^{-1}$ in notazione scientifica.
    \item Calcola e scrivi in notazione scientifica: $(7{,}5\times10^{-3})\cdot(2\times10^4)$.
    \item Calcola e scrivi in notazione scientifica: $\dfrac{3{,}6\times10^7}{1{,}2\times10^3}$.
\end{enumerate}
\begin{solution}
\begin{enumerate}
    \item $5{,}2\times10^{-4}$; $1{,}24\times10^6$
    \item $1{,}6\times10^3$
    \item $10^2=100$
    \item $\dfrac{1}{6{,}02\times10^{23}}\approx1{,}66\times10^{-24}$
    \item $1{,}5\times10^2$
    \item $3\times10^4$
\end{enumerate}
\end{solution}
\end{question}

%----------------------------------------------------%

\begin{question}[ID=1.5, subtitle={Esercizi avanzati sulle potenze}, tags={potenze, esercizi misti}, difficulty=difficile]
Risolvi o argomenta i seguenti esercizi sulle potenze:
\begin{multicols}{2}
\begin{enumerate}[leftmargin=3.2cm]
    \item Dimostra che per ogni $a\neq0$ e ogni intero $m,n$ vale $a^{m}\cdot a^{n}=a^{m+n}$ partendo dalla definizione di potenza per esponenti naturali e indicando come estendere la proprietà agli esponenti negativi.
    \item Semplifica:
    \[
    \dfrac{(2^{3}a^{-2}b)^{2}\cdot (ab^{-1})^{-3}}{2^{4}a^{-1}}
    \]
    \item Un polinomio ha termine principale $7x^{5}$. Se si moltiplica il polinomio per $x^{-3}$ e poi si eleva il risultato al quadrato, quale sarà il termine con la massima potenza di $x$?
    \item In un esperimento la concentrazione segue la legge $C(t)=C_{0}\,10^{-kt}$ con $C_{0}=2{,}5$ e $k=0{,}03$. Calcola il rapporto $\dfrac{C(10)}{C(0)}$ e determina approssimativamente dopo quante ore la concentrazione sarà metà di quella iniziale.
    \item Semplifica:
    \[
    \dfrac{(3^{-2}a^3b^{-1})^2}{3^{-3}a^2b^{-2}}
    \]
    \item Risolvi l'equazione:
    \[
    10^{2x-1}=4
    \]
\end{enumerate}
\end{multicols}
\begin{solution}
\begin{enumerate}
    \item Per gli esponenti naturali la proprietà deriva dalla somma dei fattori uguali; per gli esponenti negativi si usa $a^{-n}=\dfrac{1}{a^n}$ e si ottiene ancora $a^m\cdot a^n=a^{m+n}$.
    \item $2^2a^{-1}b^{-1}=\dfrac{4}{ab}$
    \item $49x^4$
    \item $\dfrac{C(10)}{C(0)}=10^{-0{,}3}\approx0{,}501$; il dimezzamento avviene per $t=\dfrac{\log 2}{0{,}03}\approx10{,}03$ ore.
    \item $\dfrac{a^4}{3}$
    \item $x=\dfrac{1+\log 4}{2}\approx0{,}801$
\end{enumerate}
\end{solution}
\end{question}


% ===============================

% SEZIONE 2: Espressioni frazionarie

% ===============================

\begin{question}[ID=2.1, subtitle={Calcolo di espressioni frazionarie — serie 1}, tags={espressioni, frazionarie}, difficulty=facile]
Calcola il valore delle seguenti espressioni:
\begin{multicols}{2}
\begin{enumerate}[leftmargin=3.2cm]
    \item $-\dfrac{2}{3} -\Biggl\{ \dfrac{1}{6} -\left(\dfrac{1}{2}-\dfrac{1}{3}\right) +3 -\left(\dfrac{7}{12}+\dfrac{1}{4}\right) \Biggr\}$
    \item $\left[ \left(\dfrac{5}{2}\right)^{4} \cdot \left(3-\dfrac{1}{2}\right)^{2} \right] :\left(\dfrac{5}{2}\right)^{2} :\left[ \left(-\dfrac{5}{2}\right)^{6} :\dfrac{5}{2} \right]$
    \item $\left\{ \dfrac{15}{16} - \left[ \left(\dfrac{3}{2}-\dfrac14\right)^{2} :\dfrac{5}{4} - \left(\dfrac12+\dfrac14\right)^{2} \right] \right\}^{2} +\dfrac14$
    \item $\left[ \left(0{,}20+\dfrac{2}{3}\right) : \left(\dfrac{7}{3}-2\right) \right] : \left[ 1{,}3 : \left(0{,}75-\dfrac{2}{3}\right)^{-1} \right]$
    \item $\dfrac{3}{4}-\left(\dfrac{1}{2}-\dfrac{1}{8}\right)+\dfrac{2}{3}$
    \item $\left(\dfrac{5}{3}\right)^2:\left(\dfrac{5}{3}\right)+\dfrac{1}{6}$
\end{enumerate}
\end{multicols}
\begin{solution}
\begin{enumerate}
    \item $-\dfrac{17}{6}$
    \item $\dfrac{2}{5}$
    \item $\dfrac{5}{16}$
    \item $24$
    \item $\dfrac{25}{24}$
    \item $\dfrac{11}{6}$
\end{enumerate}
\end{solution}
\end{question}

%----------------------------------------------------%

\begin{question}[ID=2.2, subtitle={Calcolo di espressioni frazionarie — serie 2}, tags={espressioni, frazionarie}, difficulty=media]
Calcola il valore delle seguenti espressioni:
\begin{multicols}{2}
\begin{enumerate}[leftmargin=3.2cm]
    \item $-\dfrac{1}{2}-\left(\dfrac{1}{4}-\dfrac{1}{2}+1\right)$
    \item $\left(\dfrac{3}{2}\right)^3:\left(\dfrac{3}{2}\right)$
    \item $\left[\dfrac{5}{6}-\left(\dfrac{1}{2}+\dfrac{1}{3}\right)\right]^2+\dfrac{1}{4}$
    \item $\left(0{,}5+\dfrac{1}{2}\right):\left(\dfrac{3}{2}-1\right)$
    \item $\left[\left(\dfrac{3}{4}+\dfrac{1}{2}\right):\dfrac{5}{3}\right]^2$
    \item $\left(0{,}6-\dfrac{1}{4}\right):\dfrac{7}{10}+\dfrac{1}{5}$
\end{enumerate}
\end{multicols}
\begin{solution}
\begin{enumerate}
    \item $-1$
    \item $\dfrac{9}{4}$
    \item $\dfrac{1}{4}$
    \item $2$
    \item $\dfrac{9}{16}$
    \item $\dfrac{7}{10}$
\end{enumerate}
\end{solution}
\end{question}


% ===============================

% SEZIONE 3: Espressioni con numeri periodici

% ===============================

\begin{question}[ID=3.1, subtitle={Espressioni con numeri periodici}, tags={espressioni, periodici}, difficulty=media]
Calcola il valore delle seguenti espressioni con numeri periodici:
\begin{multicols}{2}
\begin{enumerate}[leftmargin=3.2cm]
    \item $1{,}2-\dfrac{3}{5}\cdot2{,}\overline{8}\cdot(1{,}3\cdot10)^{-1}$
    \item $\left[0{,}2\overline{6}:(0{,}4)^2\right]^2:\left(\dfrac{5}{3}\right)^4$
    \item $0{,}\overline{3}+0{,}2\overline{6}$
    \item $1{,}\overline{2}-0{,}\overline{6}$
    \item $\left(0{,}\overline{6}-\dfrac{1}{3}\right)^2+\dfrac{1}{9}$
    \item $\left(1{,}1\overline{6}-\dfrac{1}{6}\right):0{,}\overline{6}$
\end{enumerate}
\end{multicols}
\begin{solution}
\begin{enumerate}
    \item $\dfrac{16}{15}$
    \item $\dfrac{9}{25}$
    \item $\dfrac{3}{5}$
    \item $\dfrac{5}{9}$
    \item $\dfrac{1}{9}$
    \item $1$
\end{enumerate}
\end{solution}
\end{question}